% !TeX root = thesis.tex
\section{Discussion of the findings}\label{sec:discussion}
\noindent For the receiving research team, a successful copy answers only the first question. The practical choice also depends on which measurements survive, how soon the table can be read and what happens if the connection breaks.

\subsection{What is lost when a column is removed?}
Missingness reached 65.05\% across the cohort, but it varied considerably by measurement. Heart rate was available in most hourly records, whereas lactate was missing from 96.57\%. Removing lactate would therefore remove a relatively small number of values, yet those values would be lost to any researcher interested in the occasions when lactate was measured. This is a concrete difficulty with applying a single missingness rule to a clinical record. The profiling results in Table~\ref{tab:chapter4-1} give the receiving team a basis for deciding which omissions it can accept.

At the 80\% missingness threshold, filtering alone retained 13 of 40 predictors. With the protection rules, 28 remained. The first choice saved more space, while the second preserved more measurements. The rules were exploratory and received no independent clinical endorsement. Likewise, a high correlation between two variables warranted review of what each measured before either was dropped. This extends the preprocessing question raised by \citet{Tang2020}: for a transfer, the analyst also needs a record of what the sender left out.

\subsection{Which costs matter to the receiver?}
At 16 components, the autoencoder recovered more of the variation in observed values than PCA. Its five-nearest-neighbour overlap was also higher, 0.771 against 0.436 on the original feature set. This finding is consistent with the ability of nonlinear encoders to learn compact representations described by \citet{Hinton2006}. An overall improvement can nevertheless conceal errors in individual measurements. At eight components, for instance, the autoencoder produced 54,805 newly flagged range violations against 10,179 for PCA. These are exploratory numerical screens. The ECG evaluation by \citet{Guan2026} illustrates a further test. The analyst can repeat the intended analysis on the recovered records. For this cohort, whether the reconstructions preserve performance on a patient-level task remains open.

For a research unit receiving a fresh extract, the sender's time matters. At two components, the autoencoder archive occupied 6.85 MB against 7.32 MB for PCA, yet their fresh-data delivery estimates were 649.67 s and 25.15 s respectively. Fitting alone took 609.48 s for that autoencoder. When the payloads already existed, their corresponding delivery estimates fell to 36.66 s and 21.41 s. A team choosing between those methods should first establish whether the payload will be created specifically for the delivery.

The columnar storage comparisons of \citet{Zeng2023} and \citet{Liu2025} describe trade-offs between size and processing cost. The KNUST results show how the choice of endpoint affects the practical recommendation. Parquet ranked first in the eight equal-weight, loss-tolerant fresh-data groups. Yet on the original feature set gzip had a shorter estimated fresh-data delivery time than Parquet, 27.16 s against 29.72 s. The ranking also rewarded small packages and information preservation, so a team concerned chiefly with elapsed time can reasonably make a different choice.

The direct transfers of the five Phase Two files help distinguish the effect of the representation from that of the transport. Switching from the source CSV to Parquet reduced the median copy time with either method. For Parquet itself, the median within-block difference in copy time was only 0.13 s. Across the four blocks, each method copied Parquet sooner twice. The receiving team has a clearer basis here for choosing a smaller lossless representation than for preferring one transport for every uninterrupted upload.

\subsection{What happens when an upload stops?}
The 586 successful baseline package transfers established that data could reach the university server and be recovered there. Every one completed on its first attempt. The follow-up supplied the missing recovery evidence. On the uninterrupted link, the smaller lossless archives generally reached verified publication sooner with SCP. After the break in the AE-16 upload, both methods succeeded on their first retry, but rsync finished sooner in each paired block.

The sender's copy traces help explain the difference. SCP forwarded 42.91 MB on each successful retry, approximately the size of the whole archive. Rsync forwarded between 23.07 and 23.37 MB. This pattern is consistent with the reuse of matching blocks at the receiver described by \citet{TridgellMackerras1996}. These byte counts describe the outgoing copy protocol before SSH encryption. Over the four interrupted pairs the median time saving with rsync was 10.46 s. The finding supports retaining a partial copy for a retry on a comparable route, especially when preparing the archive has already taken considerable time.

\subsection{Scope of the evidence}
The study used one public cohort and one route to KNUST. The original timing policy produced ten observations for some packages and one for others. The paired archive experiment occupied a single ten-minute window. The paired baseline files were transferred on a different day. Comparisons within each block limit the effect of changes between blocks, although the link can also change between successive transfers. Only AE-16 was deliberately interrupted, at one point in the upload. Other file sizes, outage durations, server loads and hospital links in Ghana could give different timings.

The original delivery estimates add separately measured stages and exclude initial data loading, general preprocessing, feature selection and production ingestion. The paired experiments timed the transfer of files already prepared at the sender and their verified publication. The study did not assess routine service availability or performance at other institutions. Neighbourhood comparisons used one record from each of at most 200 sampled patients. No downstream predictive task or clinician assessment established that a reduced representation was suitable for patient care.
